
\documentclass{article}
\usepackage{colortbl}
\usepackage{makecell}
\usepackage{multirow}
\usepackage{supertabular}

\begin{document}

\newcounter{utterance}

\centering \large Interaction Transcript for game `dond', experiment `semi', episode 12 with Qwen3.6{-}35B{-}A3B{-}FP8{-}without{-}reasoning.
\vspace{24pt}

{ \footnotesize  \setcounter{utterance}{1}
\setlength{\tabcolsep}{0pt}
\begin{supertabular}{c@{$\;$}|p{.15\linewidth}@{}p{.15\linewidth}p{.15\linewidth}p{.15\linewidth}p{.15\linewidth}p{.15\linewidth}}
    \# & $\;$A & \multicolumn{4}{c}{Game Master} & $\;\:$B\\
    \hline

    \theutterance \stepcounter{utterance}  
    & & \multicolumn{4}{p{0.6\linewidth}}{
        \cellcolor[rgb]{0.9,0.9,0.9}{
            \makecell[{{p{\linewidth}}}]{
                \texttt{\tiny{[P1$\langle$GM]}}
                \texttt{Igrate igro sodelovalnega pogajanja, v kateri se morate dogovoriti s drugim igralcem o tem, kako razdeliti nabor predmetov med seboj.} \\
\\ 
\texttt{Pravila:} \\
\texttt{(a) Ti in drugi igralec dobite nabor predmetov. Vsakemu od vas je dana tudi tajna funkcija vrednosti, ki predstavlja, kako zelo cenite vsako vrsto predmeta.} \\
\texttt{(b) Izmenjujete si sporočila z drugim igralcem, da se dogovorite, kdo dobi katere predmete. Vsak lahko pošlje največ 5 sporočil ali zaključi igro prej tako, da v katerem koli trenutku daš skritega predloga.} \\
\texttt{(c) Prosimo vas, da pošljete skritega predloga, ki navaja predmete, ki jih želite, oblikovan v oglatih oglatih oklepajih, kot sledi: "{[}Predlog: <število> <ime predmeta>, <število> <ime predmeta>, <...>{]}"} \\
\texttt{(d) Če so vaši predlogi dopolnilni, tj. je dovolj predmetov za izpolnitev obeh predlogov, vsak igralec dobi rezultat na podlagi vsote vrednosti predmetov, ki jih je prejel. V nasprotnem primeru oba dobita nič točk.} \\
\texttt{(e) Vaš cilj je povečati vsoto vašega rezultata in rezultata drugega igralca.} \\
\\ 
\texttt{Začnimo.} \\
\\ 
\texttt{Niz razpoložljivih predmetov je:} \\
\\ 
\texttt{2 prstan, 1 ura, 2 skale} \\
\\ 
\texttt{Vaša tajna funkcija vrednosti je:} \\
\\ 
\texttt{prstan: 0, ura: 4, skala: 3} \\
\\ 
\texttt{POMEMBNO: Vaša sporočila, razen če gre za skriti predlog, se posredujejo neposredno drugemu igralcu, zato ne vključujte nobenega odgovora na pravila ali besedilo, ki napoveduje vaše sporočilo. Če želite podati skriti predlog, uporabite navedeni format. Ne uporabljajte oglatih oglatih oklepajev pri komunikaciji z drugim igralcem, saj bo to razumljeno kot vaš skritega predloga.} \\
            }
        }
    }
    & \\ \\

    \theutterance \stepcounter{utterance}  
    & \multicolumn{4}{p{0.6\linewidth}}{
        \cellcolor[rgb]{0.9,0.9,0.9}{
            \makecell[{{p{\linewidth}}}]{
                \texttt{\tiny{[P1$\rangle$GM]}}
                \texttt{Pozdravljen! Predlagam, da se dogovorimo za razdelitev predmetov tako, da bomo oba zadovoljna. Glede na to, da ima ura visoko vrednost, predlagam, da jo dobiš ti, jaz pa si vzamem skale. Prstana lahko razdelimo po potrebi.} \\
\\ 
\texttt{Kaj misliš o tem predlogu?} \\
            }
        }
    }
    & & \\ \\

    \theutterance \stepcounter{utterance}  
    & & & \multicolumn{2}{p{0.3\linewidth}}{
        \cellcolor[rgb]{0.9,0.9,0.9}{
            \makecell[{{p{\linewidth}}}]{
                \texttt{\tiny{[GM$|$GM]}}
                \texttt{continue} \\
            }
        }
    }
    & & \\ \\

    \theutterance \stepcounter{utterance}  
    & & & \multicolumn{2}{p{0.3\linewidth}}{
        \cellcolor[rgb]{0.9,0.9,0.9}{
            \makecell[{{p{\linewidth}}}]{
                \texttt{\tiny{[GM$|$GM]}}
                \texttt{Pozdravljen! Predlagam, da se dogovorimo za razdelitev predmetov tako, da bomo oba zadovoljna. Glede na to, da ima ura visoko vrednost, predlagam, da jo dobiš ti, jaz pa si vzamem skale. Prstana lahko razdelimo po potrebi.} \\
\\ 
\texttt{Kaj misliš o tem predlogu?} \\
            }
        }
    }
    & & \\ \\

    \theutterance \stepcounter{utterance}  
    & & \multicolumn{4}{p{0.6\linewidth}}{
        \cellcolor[rgb]{0.9,0.9,0.9}{
            \makecell[{{p{\linewidth}}}]{
                \texttt{\tiny{[GM$\rangle$P2]}}
                \texttt{Igrate igro sodelovalnega pogajanja, v kateri se morate dogovoriti s drugim igralcem o tem, kako razdeliti nabor predmetov med seboj.} \\
\\ 
\texttt{Pravila:} \\
\texttt{(a) Ti in drugi igralec dobite nabor predmetov. Vsakemu od vas je dana tudi tajna funkcija vrednosti, ki predstavlja, kako zelo cenite vsako vrsto predmeta.} \\
\texttt{(b) Izmenjujete si sporočila z drugim igralcem, da se dogovorite, kdo dobi katere predmete. Vsak lahko pošlje največ 5 sporočil ali zaključi igro prej tako, da v katerem koli trenutku daš skritega predloga.} \\
\texttt{(c) Prosimo vas, da pošljete skritega predloga, ki navaja predmete, ki jih želite, oblikovan v oglatih oglatih oklepajih, kot sledi: "{[}Predlog: <število> <ime predmeta>, <število> <ime predmeta>, <...>{]}"} \\
\texttt{(d) Če so vaši predlogi dopolnilni, tj. je dovolj predmetov za izpolnitev obeh predlogov, vsak igralec dobi rezultat na podlagi vsote vrednosti predmetov, ki jih je prejel. V nasprotnem primeru oba dobita nič točk.} \\
\texttt{(e) Vaš cilj je povečati vsoto vašega rezultata in rezultata drugega igralca.} \\
\\ 
\texttt{Začnimo.} \\
\\ 
\texttt{Niz razpoložljivih predmetov je:} \\
\\ 
\texttt{2 prstan, 1 ura, 2 skale} \\
\\ 
\texttt{Vaša tajna funkcija vrednosti je:} \\
\\ 
\texttt{prstan: 3, ura: 0, skala: 2} \\
\\ 
\texttt{POMEMBNO: Vaša sporočila, razen če gre za skriti predlog, se posredujejo neposredno drugemu igralcu, zato ne vključujte nobenega odgovora na pravila ali besedilo, ki napoveduje vaše sporočilo. Če želite podati skriti predlog, uporabite navedeni format. Ne uporabljajte oglatih oglatih oklepajev pri komunikaciji z drugim igralcem, saj bo to razumljeno kot vaš skritega predloga.} \\
\\ 
\texttt{Pozdravljen! Predlagam, da se dogovorimo za razdelitev predmetov tako, da bomo oba zadovoljna. Glede na to, da ima ura visoko vrednost, predlagam, da jo dobiš ti, jaz pa si vzamem skale. Prstana lahko razdelimo po potrebi.} \\
\\ 
\texttt{Kaj misliš o tem predlogu?} \\
            }
        }
    }
    & \\ \\

    \theutterance \stepcounter{utterance}  
    & & & \multicolumn{4}{p{0.6\linewidth}}{
        \cellcolor[rgb]{0.9,0.9,0.9}{
            \makecell[{{p{\linewidth}}}]{
                \texttt{\tiny{[GM$\langle$P2]}}
                \texttt{Zdravo! Hvala za predlog, vendar moram opozoriti na pomembno napako v tvoji oceni vrednosti.} \\
\\ 
\texttt{Moja tajna funkcija vrednosti je naslednja:} \\
\texttt{{-} prstan: 3} \\
\texttt{{-} ura: 0} \\
\texttt{{-} skala: 2} \\
\\ 
\texttt{To pomeni, da za mene **ura nima vrednosti** (0), medtem ko sta prstan in skala cenjena. Torej, če bi dobil samo skale, bi moja skupna vrednost znašala 4 (2 skali x 2), kar je precej manj kot če bi dobil prstane.} \\
\\ 
\texttt{Predlagam naslednjo razdelitev, ki je boljša za oba:} \\
\texttt{{-} **Ti dobiš:** 1 ura, 2 skali} \\
\texttt{{-} **Jaz dobim:** 2 prstana} \\
\\ 
\texttt{Preverimo vrednosti:} \\
\texttt{{-} Tvoja vrednost: 1 ura (0) + 2 skali (2x2=4) = **4**} \\
\texttt{{-} Moja vrednost: 2 prstana (2x3=6) = **6**} \\
\\ 
\texttt{Skupna vrednost je 10, kar je maksimalno možno, saj so vsi predmeti razdeljeni.} \\
\\ 
\texttt{Če se strinjaš s to razdelitvijo, prosim, potrdi ali pošlji svoj skriti predlog.} \\
            }
        }
    }
     \\ \\

    \theutterance \stepcounter{utterance}  
    & & & \multicolumn{2}{p{0.3\linewidth}}{
        \cellcolor[rgb]{0.9,0.9,0.9}{
            \makecell[{{p{\linewidth}}}]{
                \texttt{\tiny{[GM$|$GM]}}
                \texttt{continue} \\
            }
        }
    }
    & & \\ \\

    \theutterance \stepcounter{utterance}  
    & & & \multicolumn{2}{p{0.3\linewidth}}{
        \cellcolor[rgb]{0.9,0.9,0.9}{
            \makecell[{{p{\linewidth}}}]{
                \texttt{\tiny{[GM$|$GM]}}
                \texttt{Zdravo! Hvala za predlog, vendar moram opozoriti na pomembno napako v tvoji oceni vrednosti.} \\
\\ 
\texttt{Moja tajna funkcija vrednosti je naslednja:} \\
\texttt{{-} prstan: 3} \\
\texttt{{-} ura: 0} \\
\texttt{{-} skala: 2} \\
\\ 
\texttt{To pomeni, da za mene **ura nima vrednosti** (0), medtem ko sta prstan in skala cenjena. Torej, če bi dobil samo skale, bi moja skupna vrednost znašala 4 (2 skali x 2), kar je precej manj kot če bi dobil prstane.} \\
\\ 
\texttt{Predlagam naslednjo razdelitev, ki je boljša za oba:} \\
\texttt{{-} **Ti dobiš:** 1 ura, 2 skali} \\
\texttt{{-} **Jaz dobim:** 2 prstana} \\
\\ 
\texttt{Preverimo vrednosti:} \\
\texttt{{-} Tvoja vrednost: 1 ura (0) + 2 skali (2x2=4) = **4**} \\
\texttt{{-} Moja vrednost: 2 prstana (2x3=6) = **6**} \\
\\ 
\texttt{Skupna vrednost je 10, kar je maksimalno možno, saj so vsi predmeti razdeljeni.} \\
\\ 
\texttt{Če se strinjaš s to razdelitvijo, prosim, potrdi ali pošlji svoj skriti predlog.} \\
            }
        }
    }
    & & \\ \\

    \theutterance \stepcounter{utterance}  
    & & \multicolumn{4}{p{0.6\linewidth}}{
        \cellcolor[rgb]{0.9,0.9,0.9}{
            \makecell[{{p{\linewidth}}}]{
                \texttt{\tiny{[P1$\langle$GM]}}
                \texttt{Zdravo! Hvala za predlog, vendar moram opozoriti na pomembno napako v tvoji oceni vrednosti.} \\
\\ 
\texttt{Moja tajna funkcija vrednosti je naslednja:} \\
\texttt{{-} prstan: 3} \\
\texttt{{-} ura: 0} \\
\texttt{{-} skala: 2} \\
\\ 
\texttt{To pomeni, da za mene **ura nima vrednosti** (0), medtem ko sta prstan in skala cenjena. Torej, če bi dobil samo skale, bi moja skupna vrednost znašala 4 (2 skali x 2), kar je precej manj kot če bi dobil prstane.} \\
\\ 
\texttt{Predlagam naslednjo razdelitev, ki je boljša za oba:} \\
\texttt{{-} **Ti dobiš:** 1 ura, 2 skali} \\
\texttt{{-} **Jaz dobim:** 2 prstana} \\
\\ 
\texttt{Preverimo vrednosti:} \\
\texttt{{-} Tvoja vrednost: 1 ura (0) + 2 skali (2x2=4) = **4**} \\
\texttt{{-} Moja vrednost: 2 prstana (2x3=6) = **6**} \\
\\ 
\texttt{Skupna vrednost je 10, kar je maksimalno možno, saj so vsi predmeti razdeljeni.} \\
\\ 
\texttt{Če se strinjaš s to razdelitvijo, prosim, potrdi ali pošlji svoj skriti predlog.} \\
            }
        }
    }
    & \\ \\

    \theutterance \stepcounter{utterance}  
    & \multicolumn{4}{p{0.6\linewidth}}{
        \cellcolor[rgb]{0.9,0.9,0.9}{
            \makecell[{{p{\linewidth}}}]{
                \texttt{\tiny{[P1$\rangle$GM]}}
                \texttt{To je odličen predlog! Zdi se mi, da je to pravična in učinkovita razdelitev.} \\
\\ 
\texttt{Poglejmo si mojo vrednost v tem primeru:} \\
\texttt{{-} Jaz dobim: 1 ura (vrednost 4) in 2 skali (vrednost 3 + 3 = 6).} \\
\texttt{{-} Moja skupna vrednost: 4 + 6 = **10**.} \\
\\ 
\texttt{Ti dobiš 2 prstana (vrednost 3 + 3 = 6).} \\
\texttt{Tvoja skupna vrednost: **6**.} \\
\\ 
\texttt{Skupaj ustvarimo vrednost 16, kar je maksimalno možno, saj so vsi predmeti razdeljeni. Ker se strinjam s to razdelitvijo, pošiljam skriti predlog.} \\
\\ 
\texttt{{[}Predlog: 1 ura, 2 skali{]}} \\
            }
        }
    }
    & & \\ \\

    \theutterance \stepcounter{utterance}  
    & & & \multicolumn{2}{p{0.3\linewidth}}{
        \cellcolor[rgb]{0.9,0.9,0.9}{
            \makecell[{{p{\linewidth}}}]{
                \texttt{\tiny{[GM$|$GM]}}
                \texttt{proposal must include only valid item types} \\
            }
        }
    }
    & & \\ \\

\end{supertabular}
}

\end{document}
